% 原文 pp.28--33、2.3節全体。
\subsectionlink{量子もつれ}
次に、デコヒーレンスの基礎となる中心的な過程、量子もつれに目を向けよう。まず量子もつれを定義する。状態ベクトル$\lvert\Psi\rangle$で記述され、二つの部分系$\mathcal S_1$、$\mathcal S_2$からなる量子系$\mathcal S$が与えられたとする（したがって$\mathcal S$は二部分量子系と呼ばれる）。$\mathcal S$の状態ベクトル$\lvert\Psi\rangle$が、これら二つの部分系の状態ベクトルのテンソル積として書けないとき、すなわち、$\lvert\Psi\rangle=\lvert\psi\rangle_1\otimes\lvert\phi\rangle_2$となるような$\mathcal S_1$の状態ベクトル$\lvert\psi\rangle_1$と$\mathcal S_2$の状態ベクトル$\lvert\phi\rangle_2$が存在しないとき、$\lvert\Psi\rangle$は$\mathcal S_1$、$\mathcal S_2$に関してもつれていると言う\footnote{記法について簡単に述べる。複数の系を扱うとき、ケットがどの系に属するかを、ケット記号の末尾の添字によって示すことが多い。たとえば「$\lvert\psi\rangle_1$」は、このケットが系$\mathcal S_1$を指すことを示す。他方、系$\mathcal A$の状態を「$\lvert a\rangle$」と書くなど、系を直接参照して量子状態に名前を付けられる場合もある。そのような場合や、同様に曖昧さのない場合には、記法を明瞭にするため添字を省略する。}。もつれた状態ではなく、もつれた系の言葉で言うと、複合系$\mathcal S$の状態が、$\mathcal S_1$、$\mathcal S_2$の何らかの状態ベクトル$\lvert\psi\rangle_1$、$\lvert\phi\rangle_2$を用いて、テンソル積$\lvert\Psi\rangle=\lvert\psi\rangle_1\otimes\lvert\phi\rangle_2$の形に表せない場合、$\mathcal S_1$と$\mathcal S_2$はもつれていると言う。本書の以後の部分では、量子状態のテンソル積について、テンソル積記号「$\otimes$」を陽には書かない。ただし、異なるHilbert空間に属する演算子の積を表すときと、Hilbert空間の積を述べるときには、記号「$\otimes$」を残す。

% 原文 p.29
\hypertarget{srcpage:29}{}\label{srcpage:29}
純粋なもつれ状態の例として、それぞれの2次元Hilbert空間の互いに直交する基底状態$\lvert0\rangle_i$、$\lvert1\rangle_i$（$i=1,2$）で記述される、二つのスピン$1/2$粒子を考えよう。状態$\lvert0\rangle_i$と$\lvert1\rangle_i$は、空間内の与えられた軸に沿って、粒子$i$のスピンがそれぞれ「上向き」と「下向き」であることに対応する。二つのスピン$1/2$粒子からなる複合系には、最大にもつれた、いくつかの異なる純粋量子状態がある（「最大の量子もつれ」の意味は後で説明する）。これらは一般にBell状態と呼ばれ、
\begin{align}
\lvert\Phi^\pm\rangle&=\frac1{\sqrt2}(\lvert0\rangle_1\lvert0\rangle_2\pm\lvert1\rangle_1\lvert1\rangle_2),\tag{2.7a}\\
\lvert\Psi^\pm\rangle&=\frac1{\sqrt2}(\lvert0\rangle_1\lvert1\rangle_2\pm\lvert1\rangle_1\lvert0\rangle_2)\tag{2.7b}
\end{align}
で与えられる。これらのBell状態を、二つのスピン$1/2$部分系の状態のテンソル積として書くことが、実際に不可能であることは、容易に直接確認できる。式(2.7)では、状態$\lvert0\rangle_1$、$\lvert1\rangle_1$が、状態$\lvert0\rangle_2$、$\lvert1\rangle_2$と一対一に相関しているため、$\lvert0\rangle_2$、$\lvert1\rangle_2$を、$\mathcal S_1$の状態$\lvert0\rangle_1$、$\lvert1\rangle_1$に対する$\mathcal S_2$の相対状態と呼ぶことが多い（逆についても同様である）。

量子もつれの定義は、何を意味しているのだろうか。$\lvert\Psi\rangle=\lvert\psi\rangle_1\lvert\phi\rangle_2$、すなわち$\mathcal S_1$と$\mathcal S_2$がもつれていない場合には、二つの部分系$\mathcal S_1$と$\mathcal S_2$を、それぞれ独立の実体と見なせる。各部分系は、その物理状態の完全な記述を構成する、自分自身の量子状態を持つ。また、複合系$\mathcal S$でのみ測定でき、個々の部分系に対する測定からは導けないような物理的性質は存在しない（そのような複合的、あるいは「大域的」な性質は、部分系間の量子的相関に具現される）。言い換えれば、部分系は、より大きな複合系の一部と見なされても、その個別性を完全に保つ。したがって全体（複合系$\mathcal S$）は、単に部分（部分系$\mathcal S_1$と$\mathcal S_2$）の和である。もちろん、これは古典物理学でなじみ深い状況である。逆に、$\lvert\Psi\rangle=\lvert\psi\rangle_1\lvert\phi\rangle_2$となるような$\mathcal S_1$、$\mathcal S_2$の部分系状態$\lvert\psi\rangle_1$、$\lvert\phi\rangle_2$が存在しないなら、上述の分離可能性の「古典的」特徴はすべて成り立たなくなる。この場合、部分系$\mathcal S_1$と$\mathcal S_2$に各自の量子状態を帰属させることはできず、大域的な複合量子状態によってのみ記述できる（図1.2も参照）。

「量子もつれ」（英語のentanglement、ドイツ語のVerschränkung）という語は、1935年にSchrödingerが初めて作った\cite{ref002,ref052,ref053}。彼はすぐに、その非古典的な含意を強調した\cite[p.~555]{ref052}。
\begin{quote}
それぞれを表すものによって状態が分かっている二つの系が、両者の間の既知の力により一時的な物理的相互作用に入り、相互に影響し合う期間を経て再び分離すると、それらを以前と同じように、各自にそれぞれを表すものを与えることによって記述することは、もはやできない。私はこれを、量子力学の一つの特徴というよりも、量子力学を特徴付けるまさにその性質、古典的な思考からの全面的な離脱を強いる性質と呼びたい。
\end{quote}

% 原文 p.30
\hypertarget{srcpage:30}{}\label{srcpage:30}
量子もつれの特異な特徴は、現在までに多数の実験で確かめられている。たとえば、何キロメートルも離れた、もつれた光子対を用いる実験である\footnote{量子もつれ、Bell不等式、波動関数の収縮理論に関する原理実証実験のため、こうした実証を宇宙空間へ拡張して数千キロメートルの隔たりを達成する展望については、文献\cite{ref054}を参照。}。量子もつれは、量子コンピュータや量子暗号などの量子技術に関わる「第二次量子革命」\cite{ref055}の核心にもある。量子コンピュータは、ある種の計算問題を、どの古典コンピュータにも決して不可能な速さで解けるはずであり（第7章）、量子暗号は完全に安全な通信を可能にする。相当に複雑な量子コンピュータを実験的に実現できるか、またそれがいつになるかは今後の課題であるが、量子暗号はすでに比較的成熟した分野であり、商業的な応用も存在する。

\subsubsectionlink{量子的相関と古典的相関}
量子的相関、すなわち量子もつれは、古典的相関と根本的に異なることに注意するのが重要である\footnote{Schrödinger自身も、量子もつれについての初期の論文で、この相違を常に十分明確にしたわけではないように見える。たとえば二つの論文には「分離した系の間の確率関係についての考察」という共通の題名が付いている\cite{ref052,ref053}が、これは、古典的な確率関係をはるかに超える、量子もつれの特異な量子的性格を反映していない。}。古典物理学では、相関はしばしば、ある保存則によって生じる。たとえば静止した粒子が、同一の二つの破片に崩壊すると、全運動量の保存により、破片は同じ速さで反対方向へ飛び去る。一方の破片の運動量を測れば、他方の運動量は、大きさが等しく符号が反対でなければならないと、直ちに推論できる。これは二粒子間の純粋に古典的な、統計的な一対一の相関である。各破片の運動量は、第一の破片に行う測定とは独立に「存在」し、第二の粒子の運動量の推論は、運動量保存から直接に従う。

次に量子もつれの場合に移ろう。スピン$1/2$粒子の対が、Bell状態［式(2.7b)］
\begin{equation}\lvert\Psi^+\rangle=\frac1{\sqrt2}(\lvert0\rangle_1\lvert1\rangle_2+\lvert1\rangle_1\lvert0\rangle_2)\tag{2.8}\end{equation}
で記述されるとする。第一の粒子を測定して「上向きスピン」が見つかった場合、第二の粒子を測定すれば必ず「下向きスピン」が見つかると直ちに推論できる、という範囲では、古典的相関との類似が成り立つ\footnote{最初のEPR論文\cite{ref010}の有名な表現を用いれば、系2の測定結果を「確実に予言」できる。したがって直観的には、第二の系は、この系への明示的な測定によって予言が確かめられるより前から、すでに「下向き」状態にあったと結論するだろう。本質的に、この見方はEPRの、よくも悪くも有名な「実在性の基準」に対応している。}。

% 原文 p.31
\hypertarget{srcpage:31}{}\label{srcpage:31}
そのため、最初の測定という行為によって、量子的相関が古典的な、純粋に統計的な相関へと変換されると言える。ただし量子的な設定では、第一の粒子の測定結果は完全に偶然的である。「上向き」と「下向き」は等しい確率で得られるが、具体的にどちらの結果になるかを予言する方法はない。したがって第一の粒子の測定結果に関する情報が、任意の距離だけ空間的に隔たっていてもよい第二の粒子へ、見かけ上瞬間的に伝わるように見える（この観察は、EPRの「パラドックス」\cite{ref010}の基礎である）。

一見すると、これは、どんな古典的信号も光速を超えて伝わることはできないという特殊相対論の原理に反するように思われる。しかし実際には、そのような違反はない。第二の粒子の測定結果は、第一の粒子の測定結果によって瞬間的に固定されるが、後者が完全に偶然的であるため、二者間で有用な情報を伝えることはできない。この「信号伝達不可能定理」は、二つの部分系に対する測定の統計を用いて、より厳密に証明できる（たとえば文献\cite{ref056}）。これは、かなり注目すべき結果である。量子力学それ自体は非相対論的な理論で、情報の伝達速度に根本的制限を導入するような公理を陽には含まない。言い換えれば、量子力学は特殊相対論を「知らない」。それにもかかわらず、その確率構造は間接的に、上に述べたような信号伝達不可能という制約を課す。

このような見かけ上瞬間的な「不気味な遠隔作用」（Einsteinの表現）を考えるよい方法は、もつれ状態が、過去のある時点で二つの部分系を相互作用させることによって、局所的に準備されたはずだと理解することである。その後、二つの部分系を引き離すと、量子状態は単に、より広い空間領域に「広がる」、すなわち非局在化する。二つの部分系の間で物理的に伝達される必要のある情報はない。もつれ状態は、すでに個々の部分系についての最も完全な記述であり、これらの系についての可能なすべての情報を包含している。

この見方は、Schrödingerの1935年の最初期の量子もつれ論文\cite{ref002}にも現れている（J.~D.~Trimmerによる英訳\cite{ref057}から引用）。
% 原文 p.32
\hypertarget{srcpage:32}{}\label{srcpage:32}
\begin{quote}
知識の一部が、選言的な条件文の形で二つの系の間に漂っているということは、世界の反対側から二つを持ってきて、相互作用なしに並べた場合には、確かに起こりえない。その場合、両者は実際、互いについて何も「知らない」からである。一方への測定によって、他方に何を予想すべきかを把握できるはずはない。生じる「予言のもつれ」は、二つの物体が以前のある時に、真の意味で一つの系を形成していたこと、すなわち相互作用しており、互いに痕跡を残したことにしか、由来しえないのは明らかである。
\end{quote}

量子もつれという性質は、自然が根本的に非局所的であることを示唆するように見える。第二の粒子に対する局所的な測定の結果は、複合系の大域的なもつれた量子状態にだけ符号化されている、量子的相関によって決まるのである。

\subsubsectionlink{量子もつれと識別可能性の定量化}
四つのBell状態(2.7)を導入した際、これらが「最大」にもつれていると述べた。これは、量子もつれに異なる定量的程度が存在することを示唆する。そこで、量子もつれをどのように定量化するか、つまり、与えられた状態が「どれほど」もつれているかを、どう測るかという問いが生じる。この問いはデコヒーレンスにとって重要である。後に見るように、大まかに言えば、関心のある系と環境の間のもつれの程度が大きいほど、通常、環境のデコヒーレンス効果は強くなる。

次の形の、二部分のもつれ状態を考えよう。
\begin{equation}\lvert\Psi\rangle=\frac1{\sqrt2}(\lvert\psi_1\rangle_1\lvert\phi_1\rangle_2\pm\lvert\psi_2\rangle_1\lvert\phi_2\rangle_2).\tag{2.9}\end{equation}
ここで$\lvert\psi_i\rangle_1$、$\lvert\phi_i\rangle_2$（$i=1,2$）は、部分系$\mathcal S_1$、$\mathcal S_2$の任意の状態で、必ずしも互いに直交していない。この状態にある量子もつれを定量化する、直観的で有用な方法は、次の問いを考えることである。一方の系を測定することによって、観測者は他方について、どれだけ知ることができるだろうか。

Bell状態(2.7)を振り返ろう。これらが最大にもつれているのは、系$\mathcal S_2$を$\{\lvert0\rangle_2,\lvert1\rangle_2\}$基底で射影測定すれば、その後の測定で系$\mathcal S_1$を$\lvert0\rangle_1$、$\lvert1\rangle_1$のどちらに見いだすかが、直ちに分かるからである。この事実は二つの性質による。第一に、$\mathcal S_2$の「相対」状態$\lvert0\rangle_2$、$\lvert1\rangle_2$が、$\mathcal S_1$の状態$\lvert0\rangle_1$、$\lvert1\rangle_1$と一対一に相関している。第二に、状態$\lvert0\rangle_2$、$\lvert1\rangle_2$が互いに直交しており、完全に識別できる。これらの状態を、系$\mathcal S_1$の相対状態を示す、装置の何らかの指針に対応すると考えてよい。実際、これはまさに、後の2.5.1節で論じるvon Neumannの測定方式の基本的な考えである。

% 原文 p.33
\hypertarget{srcpage:33}{}\label{srcpage:33}
次に状態(2.9)について、$\mathcal S_2$の状態$\lvert\phi_1\rangle_2$、$\lvert\phi_2\rangle_2$が直交から遠く離れ、大きく重なり合う極限を考えよう。すると、$\mathcal S_2$の射影測定でこの二つの状態を識別するのは、とても難しい。その結果、この測定から$\mathcal S_1$の対応する相対状態（$\lvert\psi_1\rangle_1$か$\lvert\psi_2\rangle_1$か）を推論するのも難しい。大まかに言えば、系$\mathcal S_2$は、状態$\lvert\psi_1\rangle_1$と$\lvert\psi_2\rangle_1$に関して$\mathcal S_1$を区別する「情報」を、ほとんど符号化していないと言える。

状態$\lvert\phi_1\rangle_2$、$\lvert\phi_2\rangle_2$が等しいという極端な場合、$\lvert\phi_1\rangle_2=\lvert\phi_2\rangle_2\equiv\lvert\phi\rangle_2$なので、式(2.9)を積状態
\begin{equation}\lvert\Psi\rangle=\frac1{\sqrt2}(\lvert\psi_1\rangle_1\pm\lvert\psi_2\rangle_1)\lvert\phi\rangle_2\tag{2.10}\end{equation}
と書ける。したがって、この状態はもはやもつれておらず、両系間の量子的相関を含まない。いま、どんな基底でもよいから系$\mathcal S_2$を測定しても、その後$\mathcal S_1$を射影測定したときに、$\lvert\psi_1\rangle_1$と$\lvert\psi_2\rangle_1$のどちらが見つかるかについては、何も推論できない。二つの部分系はそれぞれ別々の純粋状態にあり、その各状態が、個々の系の物理状態を完全に指定する。より認識論的に言えば、たとえば測定統計に関して「知りうるすべて」を指定する。「量子的な全体性」という特徴は消失している。最後に、上の議論は、量子もつれの程度と、$\mathcal S_1$の状態$\lvert\psi_1\rangle_1$、$\lvert\psi_2\rangle_1$の重なりとの相関についても、同じように述べられる。

間もなく見るように、他の系の状態と相関している、一方の系の状態のこの識別可能性は、デコヒーレンスの概念的基礎を理解する核心にある。ここでは部分系$\mathcal S_2$は、式(2.9)の形の量子的相関を介して、$\mathcal S_1$に関する「情報」を符号化する環境に対応する。上の議論によれば、この情報量は、系と環境のもつれが増すとともに増加し、したがって、系の異なる成分状態と相関した環境の相対状態の識別可能性とともにも増加する。環境が獲得した系についての情報量が大きくなるほど、系は、個別の量子状態を割り当てられないという意味で、その個別性を失う。その結果、初めは系の内部に局在していた量子コヒーレンスは、系と環境の複合状態の「共有された性質」となり、系の水準ではもはや観測できなくなって、デコヒーレンスに至る。以下の2.6節と2.7節で、これらの考えをより正確にする。
